\chapter{Geometric Foundations}
\label{chapter:Geometric-Foundations}

%&& a &= b & \text{Proof \ref{proof:}} \label{equation:}

\section{Topology}

euclidean space relate to RN, essentially affine space

homeomorphism

\section{Manifolds}

starting with set of locations, and some way to measure distances and angles

not establishing what these are but trusting they are enough for continuity

explain real and finite dimension

\begin{definition}[Manifold]
	A manifold $M$ is a second-countable Hausdorff space that is locally homeomorphic to a Euclidean space.
\end{definition}

This local homeomorphism is established with maps to $\mathbb{R}^{n}$, called charts. This is a way of assigning local coordinate systems. A collection of charts is fittingly called an atlas, and maps out every location on the manifold to at least one set of coordinates.

As the names suggest, charts and atlases may not represent the actual geometry of the manifold. Opening up a world atlas, projections like equirectangular and Mercator warp the roughly spherical surface of the geoid, making sizes inconsistent, straight lines appear curved, and curved lines appear straight. Generally, norms and inner products on $\mathbb{R}^{n}$ do not capture the geometry of a manifold. Instead, other objects will be introduced.

\begin{definition}[Chart]
	A chart $\phi: U \leftrightarrow \phi(U) \subseteq \mathbb{R}^{n}$ is a homeomorphism from an open subset of a manifold $U \subseteq M$ to a subset of real numbers $\mathbb{R}^{n}$.
\end{definition}

A chart can be split to extract individual coordinate functions $x^{i}: U \rightarrow \mathbb{R}$, with the convention $\phi(p) = (x^{1}(p), \dots, x^{n}(p))$.

\begin{definition}[Atlas]
	An atlas is is a collection of charts whose domains cover a whole manifold.
\end{definition}

\begin{definition}[Transition Map]
	A transition map $\phi_{b} \circ \phi_{a}^{-1}: \mathbb{R}^{n} \supseteq \phi_{a}(U_{a} \cap U_{b}) \leftrightarrow \phi_{b}(U_{a} \cap U_{b}) \subseteq \mathbb{R}^{n}$ is a continuous map between the coordinates of overlapping charts.
\end{definition}

\section{Vector Spaces on a Manifold}

Bla bla bla first we start with tangent space, set of all velocities/directions, geometric intuition, now definition using derivatives since derivatives correspond to velocity.

\subsection{Tangent Space}

\begin{definition}[Tangent Space]
	The tangent space $T_{p}M$ of a point on a manifold $p \in M$ is a vector space representing all of the directions one can travel at that location. It may be constructed as $T_{p}M = C_{p} / \sim$, where $C_{p}$ is the set of all continuous and differentiable curves passing through $p$, and $\sim$ is an equivalence relation determining whether two curves are tangent at $p$.
\end{definition}

The big challenge to overcome with this definition is that $M$ is just a topological space. Continuity is possible, but the space may not even have a metric! What does it mean for curves to be differentiable? The solution is to represent the curves as functions from reals to reals, where continuity and differentiability are easy to define. They are parameterized with an input argument in $\mathbb{R}$, and mapped with a chart to $\mathbb{R}^{n}$. Because transition maps are constrained to be smooth, if these are shown for one chart $(U, \phi)$, they apply for all relevant charts in the neighborhood.

Continuous and differentiable curves are ones which are continuous and differentiable when composed with any chart. The domain of $(-1, 1)$ and fixing $p$ at 0 are just parameterization conventions for simplicity.

\begin{equation}
	C_{p} = \{ \gamma : (-1, 1) \rightarrow U : \gamma(0) = p \land (\phi \circ \gamma) \in C^{1} ((-1, 1), \mathbb{R}^{n}) \} \\
\end{equation}

In the same way, tangent curves are ones which are tangent through any chart.

\begin{equation}
	\gamma_{1} \sim \gamma_{2} \iff \dv{t}(\phi \circ \gamma_{1}) (0) = \dv{t}(\phi \circ \gamma_{2}) (0)
\end{equation}

The quotient set $C_{p} / \sim$ then contains equivalence classes $[ \gamma ]_{\sim}$ corresponding to the unique tangents at $p$. These tangents form a vector space isomorphic to $\mathbb{R}^{n}$ itself (Proof ??).

A basis can be constructed for the tangent space, with respect to a particular chart $\phi$, by choosing curves $\gamma_{i}$ which vary in only one coordinate. These can be constructed backwards from appropriately restricted $c_{i}$t.

\begin{align*}
	c_{i}(t) &= (p^{1}, \dots, p^{i} + t, \dots, p^{n}) : \mathbb{R} \rightarrow \mathbb{R}^{n} \\
	\gamma_{i} &= \phi^{-1} \circ c_{i} \\
	\vb{e}_{i}& = [\gamma_{i}]_{\sim}
\end{align*}

Intuitively, these vectors point along the coordinate lines intersecting at $p$.

\begin{definition}[Action of a Vector on a Function]
	Any vector $\vb{v} = [\gamma]_{\sim} \in T_{p} M$ in the tangent space at $p$ may act on a function $f: U \rightarrow \mathbb{R}$, defined in an open neighborhood $U \subseteq M$ with the following definition.
	\begin{equation}
		\vb{v}(f) = \dv{t} (f \circ \gamma)(0)
	\end{equation}
\end{definition}

This is a consistent definition because every vector corresponds to a unique equivalence class of curves, which all have the same derivative at $p$. This suggests a correspondence to directional derivatives, and indeed the basis vectors can be identified with derivative operators.

\begin{align*}
	\vb{e}_{i}(f) &= \dv{t} (f \circ \gamma_{i})(0) \\
	&= \dv{t} (f \circ \phi^{-1} \circ c_{i})(0) \\
	&= \sum_{j = 1}^{n} \pdv{x^{j}} (f \circ \phi^{-1}) (c_{i}(0)) \cdot [c_{i}'(0)]_{j} \\
	&= \pdv{x^{i}} (f \circ \phi^{-1}) (p^{1}, \dots, p^{n})
\end{align*}

Using looser notation, $\vb{e}_{i}(f) = \pdv{f}{x^{i}}$ at $p$. Note that $[c_{i}'(0)]_{j}$ refers to the $j$th component of $c_{i}'$ at $t = 0$, which is 0 if $j \neq i$ and 1 if $j = i$.

\subsection{Cotangent Space}

\begin{definition}
	aaa
\end{definition}

\section{Vector Fields on a Manifold}

For any curve in the equivalence class, 

Now that a vector space is defined at every point on the manifold, often we want to assign a particular vector at every location. For example, everywhere on the globe, the wind blows in a particular direction. This is called a vector field.

\begin{definition}[Tangent Bundle]
	The tangent bundle $TM$ of a manifold $M$ is the disjoint union of every tangent space on the manifold.
	\begin{equation}
		TM = \bigsqcup_{p \in M} T_{p} M = \bigcup \{ (p, \vb{v}) : p \in M, \vb{v} \in T_{p} M \}
	\end{equation}
\end{definition}

A disjoint union merges sets while distinguising which set each element originated from. bla bla bla reason

\begin{definition}[Vector Field]
	A vector field $\mathbf{u} : M \hookrightarrow TM$ maps every location $p$ on a manifold $M$ to a vector in that location's tangent space $T_{p} M$.
\end{definition}

By convention, the location in the ordered pair is redundant and ignored, so $\mathbf{u}(p) = (p, \mathbf{u}_{p})$ is understood to refer just to $\mathbf{u}_{p}$.

\begin{definition}[Differential]
	content...
\end{definition}

\begin{definition}[Cotangent Space]
	content...
\end{definition}

\begin{definition}[Cotangent Bundle]
	content...
\end{definition}

vector fields as vector space w scalar at every location

vector spaces are over what field

something something coordinates

basis field

vector field that is also a vector space

notation overloading

derivative on manifold assumes using appropriate coordinates

tangent bundle

special disjoint union to maintain that every vector space is tied to a location

use globe as example

\begin{definition}[Cotangent Space]
	The cotangent space $T_{p}^{*} M$ is the dual of the tangent space at the same point $(T_{p} M)^{*}$.
\end{definition}

explain differential forms, exterior derivative, dual basis link here

\section{Introducing Geometry to a Manifold}

bla bla bla coordinate system is no good for measuring distances, e.g. lat lon

\begin{definition}[Metric Tensor]
	A metric tensor $g$ assigns a function $g(p) = g_{p}: T_{p} M \times T_{p} M \rightarrow \mathbb{R}$ to every point $p$ on a manifold with the following properties.
	\begin{enumerate}
		\item Bilinear
		\item Symmetric
		\item Non-degenerate
	\end{enumerate}
	This bilinear form varies continuously across the manifold.
\end{definition}

Because of the isomorphism between bilinear forms and tensors, $g_{p} \in T_{p}^{*} M \otimes T_{p}^{*} M$. It is a symmetric tensor of type $(0, 2)$. The components of the tensor using a particular basis are $g_{ij} = g(\vb{e}_{i}, \vb{e}_{j})$.

The introduction of a length functional finally brings the concept of distances to a manifold. The length of a curve $\gamma$ is defined as follows.

metric space

inverse metric tensor??

\begin{equation}
	L[\gamma] = \int_{D(\gamma)} \sqrt{g_{\gamma(s)}\left(\dv{t} \gamma(s), \dv{t} \gamma(s)\right)} \dd{t}
\end{equation}

\begin{definition}[Metric Signature]
	The metric signature $(v, p, r)$ of a metric tensor $g$ is the number of $v$ positive, $p$ negative, and $r$ zero eigenvalues at a point $p \in M$.
\end{definition}

The non-degeneracy requirement means that $r = 0$, and the signature is constant throughout the manifold.

show pointwise vs field

signature

\begin{definition}[Riemannian Metric]
	A Riemannian metric is a metric tensor that is additionally positive-definite, i.e. an inner product.
\end{definition}

\begin{definition}[Riemannian Manifold]
A Riemannian manifold is a sufficiently smooth manifold with a Riemannian metric.
\end{definition}

The introduction of an inner product finally brings the concept of norms and angles to the manifold, based on the law of cosines.

dot product here??

\begin{align}
	\norm{\vb{u}} &= \sqrt{g_{p}(\vb{u}, \vb{u})} \\
	\cos(\theta) &= \frac{g_{p} (\vb{u}, \vb{v})}{\norm{\vb{u}} \norm{\vb{v}}}
\end{align}

\begin{definition}[Christoffel Symbols]
	content...
\end{definition}

\begin{align*}
	\Gamma^{i}{}_{jk} = \left\langle \vb{e}^{i*}, \pdv{\vb{e}_{j}}{x^{k}} \right\rangle
\end{align*}

\section{Exterior Geometry}

Wedge product

Relate to cross product and skew maps

Never combine a vector and a normal!

Symmetry of the universe or whatever

Exterior derivative

Determinant of tensor

\section{Checklist}

topology

manifolds

(pseudo) riemannian

charts, atlases, coordinate maps

coordinate systems, several examples for different geometries

derivatives?

euclidean, elliptic, hyperbolic geometries

minkowski and lorentz spaces

curvature

using metric to define curvature

tangent space

tangent bundle

constructions

directional derivatives

frames of reference, galilean, lorentzian transformations

cartan 2 forms against christ awful

complexification in another chapter

differentials and derivatives as bases, dual pairs

parallel transport